% ============================================================================
% DISCUSSION --- three paragraphs: synthesis -> limitations -> implication.
% Compressed from an earlier four-paragraph version (archived) in which
% paragraphs 1-2 made the same move twice.
% ============================================================================
\section{Discussion and limitations}
\label{sec:discussion}

Taken together, our results support a view of molecular design as control of a learned executable process rather than endpoint generation followed by scoring or repair. \method{} separates which molecular transitions are \emph{possible}, which legal transitions are \emph{plausible}, and which plausible futures serve the current \emph{purpose}. The learned reference law captures state-dependent transition preferences beyond executable support alone, while finite-horizon control can redirect this frozen law according to future reachability. Because every committed state is itself a molecule, control can also be reapplied to states already reached and constraints can act on the trajectory rather than only its endpoint. Variable-size editing, stochastic molecular processes, and $h$-transforms are not individually new; the distinction is that \method{} brings them together in a common molecular transition process whose executable support and learned plausibility remain separate from task-specific purpose.

This decomposition also makes the limitations explicit. Control cannot create a useful continuation that is absent from the executable support or assigned negligible probability by the reference process. Exact finite-horizon guarantees apply under exact backward values, whereas molecular-scale control uses an approximate learned $h_\phi$; improved task performance therefore demonstrates useful future-aware steering, not exact conditioning at scale. Property-directed control additionally inherits errors and extrapolation from the property models defining desirability. Finally, the declared state space remains a 2D molecular representation: connectedness and valence admissibility do not imply stereochemical correctness, conformational plausibility, synthesizability, stability, or biological activity. These properties require richer representations or independent downstream evaluation.

The natural next step is therefore to improve the \emph{molecular process itself}, not only the controller placed on top of it. Richer executable state spaces and reference models could expand which chemically meaningful futures are reachable, while new measurements could update which of those futures are preferred without redefining the underlying transition process. More broadly, \method{} suggests that when a domain admits meaningful executable transformations between valid states, those transformations can form the native event language of a learned stochastic process, with generation and control acting on the same evolving object.
